\section{Correctness of the Higher-Order Interpretation}
\label{sec:definability}

The interpretation of the higher-order language at higher types mixes
the interpretation of the underlying computation with the
interpretation of the semiring dependency matrices
(c.f. \thmref{fam-closed}). It is therefore not immediate that
intepreting a program in $\Fam(\SemiMod(S))$ will give us the expected
interpretation purely on first order data. Note that the fullness of
the functor $H : \Fam(\SDSemiMod(S)) \to \Fam(\SemiMod(S))$ does mean
that we are guaranteed that the tangent maps are correctly computed.
\citet{vakar22} and \citet{nunes2023} construct custom
instances of categorical sconing arguments to prove correctness of
their higher-order interpretation with respect to normal
differentiation. Instead of doing this, we make use of a general
\emph{syntax free} theorem due to \citet{fiore-simpson99}. The proof
of this depends on the construction of a Grothendieck Logical Relation
over the extensive topology on the category $\cat{C}$, but the
statement of the theorem does not rely on this. We have formalised
this proof in Agda (see \texttt{conservativity.agda} in the
supplementary material).

\begin{theorem}[\citet{fiore-simpson99}]
  \label{thm:glr-definability}
  Let $\cat{C}$ be an extensive bicartesian category, $\cat{D}$ be a
  bicartesian closed category, and $F : \cat{C} \to \cat{D}$ a functor
  preserving finite products and coproducts. Then there is a category
  $\GLR(F)$ and functors $p : \GLR(F) \to \cat{D}$ and
  $\hat{F} : \cat{C} \to \GLR(F)$, such that:
  \begin{enumerate}
  \item $\GLR(\cat{D},F)$ is bicartesian closed;
  \item $F = p \circ \hat{F} : \cat{C} \to \cat{D}$;
  \item The functor $p$ strictly preserves the bicartesian closed structure; and
  \item The functor $\hat{F}$ is full and preserves the bicartesian structure.
  \end{enumerate}
\end{theorem}

\begin{remark}
  Compared to the exact result stated at the end of
  \citet{fiore-simpson99}'s paper, we have made two modifications,
  justified by our Agda proof. First, we generalise to the case where
  $\cat{C}$ is not Cartesian closed, and the functor $F$ does not
  preserve exponentials. Examination of the proof reveals that if this
  is the case, then $\hat{F}$ also preserves exponentials, but it is
  not needed for the result stated. Second, Fiore and Simpson restrict
  to the case when $\cat{C}$ is small to be able to construct
  Grothendieck sheaves on this category. We use Agda's universe
  hierarchy to simply construct ``large'' sheaves at the the
  appropriate universe level.
\end{remark}

We use \thmref{glr-definability} to show that the
interpretation of the language in the category $\Set$ agrees with the
higher-order interpretation in $\Fam(\SemiMod(S))$
on the underlying function at first order, as required. This shows that the
higher-order interpretation does what we expect in the underlying
interpretation of terms, and that the approximation information does
not interfere.

\begin{theorem}
  \label{thm:underlying-interp-equal} For all
  $\Gamma \vdash M : \tau$, where $\Gamma$ and $\tau$ are first-order,
  the underlying function in the interpretation
  $\sem{\Gamma \vdash M : \tau}_{\Fam(\SemiMod(S))}$ is equal to the
  interpretation $\sem{\Gamma \vdash M : \tau}_{\Set}$ in $\Set$.
\end{theorem}

\begin{proof}
  Instantiate \thmref{glr-definability} with the functor
  \begin{displaymath}
    \langle \mathrm{Id} , \pi_1 \rangle : \Fam(\SemiMod(S)) \to \Fam(\SemiMod(S)) \times \Set
  \end{displaymath}
  that is the identity in the first component and projects out the
  underlying function in the second. By
  \lemref{pi1-preserve-products-and-coproducts}, this functor
  preserves products and coproducts. For each
  $\Gamma \vdash M : \tau$, we obtain a $g$ such that
  $g = \sem{\Gamma \vdash M : \tau}_{\Fam(\SemiMod(S))}$ and
  $\pi_1(g) = \sem{\Gamma \vdash M : \tau}_\Set$. Substituting $g$
  yields the result.
\end{proof}

